\BourbakiExercises{Exercises}

\begin{enumerate}
\def\labelenumi{\arabic{enumi})}
\tightlist
\item
  Let \(\mathrm{K}\) be a commutative field. Let \(\mathrm{A}\) be the K-algebra with basis consisting of the unit element and two elements \(e, f\) such that \(e^2 = e\) , \(ef = f\) , and \(fe = f^2 = 0\) . a) Prove that the regular and coregular representations of \(\mathrm{A}\) are not isomorphic (calculate \(\operatorname{Tr}_{\mathrm{A / K}}(e)\) and \(\operatorname{Tr}_{\mathrm{A}^{\circ} / \mathrm{K}}(e)\) ).
\end{enumerate}

\begin{enumerate}
\def\labelenumi{\alph{enumi})}
\setcounter{enumi}{1}
\tightlist
\item
  Prove that the regular representation contains a subrepresentation that is isomorphic to neither a subrepresentation nor a quotient representation of the coregular representation. Moreover, there exists a simple representation of A that is not isomorphic to a subrepresentation of the regular representation.
\end{enumerate}

\begin{enumerate}
\def\labelenumi{\arabic{enumi})}
\setcounter{enumi}{1}
\tightlist
\item
  Let \(A\) be a K-algebra. We view the dual \(A^*\) of \(A\) as a left \(A\) -module.
\end{enumerate}

\begin{enumerate}
\def\labelenumi{\alph{enumi})}
\item
  The mapping \(E \mapsto E'\) is a surjection from the set of A-submodules of \(A^{*}\) to the set of right ideals of A. When E is such a submodule, the orthogonal of \(E'\) in \(A^{*}\) is a submodule of \(A^{*}\) containing E.
\item
  Prove that the orthogonal in A of the socle of \(\mathbf{A}^*\) is the radical of A.
\item
  Prove that the orthogonal in \(A^{*}\) of the socle of \(A_{d}\) is the radical of the A-module \(A^{*}\) . Deduce that \(A^{*}\) is a semisimple A-module if and only if the ring A is semisimple.
\end{enumerate}

¶ 3) Let A be an Artinian ring and \(\overline{\mathrm{A}}\) the ring \(\mathrm{A} / \Re (\mathrm{A})\) . We take a decomposition of A as a direct sum of indecomposable ideals \(\mathrm{Ae}_i\) ( \(i\in \mathrm{I}\) ), where the \(e_i\) are orthogonal idempotents (VIII, p.~179, Exercise 6, d)). We denote by \(\mathrm{I} = \cup_{\alpha \in \mathrm{R}}\mathrm{I}_{\alpha}\) the partition of I such that two elements \(i,j\) of I belong to the same set \(\mathrm{I}_{\alpha}\) if and only if the modules \(\mathrm{Ae}_i\) and \(\mathrm{Ae}_j\) are isomorphic. The \(\overline{\mathrm{A}}\) -modules \(\bigoplus_{i\in \mathrm{I}_{\alpha}}\overline{\mathrm{A}}\overline{e}_i\) are the isotypical components of \(\overline{\mathrm{A}}\) (VIII, p.~179, Exercises 6, d) and 7, a))

\begin{enumerate}
\def\labelenumi{\alph{enumi})}
\item
  Prove that if the ring A is autoinjective (VIII, p.~53, Exercise 8), then there exists a permutation \(\sigma\) of R such that every right ideal \(e_i\mathrm{A}\) , for \(i\in \mathrm{I}_{\alpha}\) , contains a unique minimal right ideal isomorphic to \(\overline{e}_j\overline{\mathrm{A}}\) for \(j\in \mathrm{I}_{\sigma (\alpha)}\) (observe that the right annihilator of an idempotent \(e\) is \((1 - e)\mathrm{A}\) ). Deduce that the left socle \(\mathfrak{s}\) of A is contained in its right socle \(\mathfrak{t}\) (prove that every right ideal \(e_i\mathfrak{s}\) is minimal, by proving that it is the annihilator of a maximal left ideal, cf.~VIII, p.~178, Exercise 2). Conclude that we have \(\mathfrak{s} = \mathfrak{t}\) and that every left ideal \(\mathrm{Ae}_i\) , for \(i\in \mathrm{I}_{\sigma (\alpha)}\) , contains a unique minimal left ideal isomorphic to \(\overline{\mathrm{A}}\overline{e}_j\) for every \(j\in \mathrm{I}_{\alpha}\) .
\item
  Conversely, if every left ideal \(Ae_{i}\) contains a unique minimal left ideal and every right ideal \(e_{i}A\) contains a unique minimal right ideal, then the ring A is autoinjective. (Use Exercise 7 of VIII, p.~179 to prove that every foot of s is a minimal two-sided ideal. Deduce that t contains s, and then that s = t. Finally, prove that condition \(\beta\) ) of Exercise 8 of VIII, p.~53 is satisfied.)
\item
  Suppose that A is an algebra of finite degree over a commutative field and an autoinjective ring. Let \(\alpha \in R\) , \(i \in I_{\alpha}\) , and \(j \in I_{\sigma(\alpha)}\) . Prove that the linear representation of A associated with the A-module \(Ae_{j}\) is isomorphic to the transpose of the linear representation of \(A^{\circ}\) associated with the right A-module \(e_{i}A\) . (Let \(U_{i}\) be the left submodule of \(A^{*}\) that is the orthogonal of \((1 - e_{i})A\) . Deduce from Exercise 7, d) of VIII, p.~179 that \(U_{i}\) is isomorphic to a quotient module of \(Ae_{j}\) and that we consequently have long \(e_{i}A \leqslant long e_{j}A\) . Conclude that these lengths are equal, so \(U_{i}\) is isomorphic to \(Ae_{j}\) . Use b) and Exercise 2, a) to prove a converse of this statement.
\end{enumerate}

\begin{enumerate}
\def\labelenumi{\arabic{enumi})}
\setcounter{enumi}{3}
\tightlist
\item
  Keep the assumptions of Exercise 3.
\end{enumerate}

\begin{enumerate}
\def\labelenumi{\alph{enumi})}
\item
  Prove that if A is a Frobenius ring (VIII, p.~53, Exercise 9), then we have \(\operatorname{Card}(\mathrm{I}_{\alpha}) = \operatorname{Card}(\mathrm{I}_{\sigma(\alpha)})\) for every \(\alpha \in R\) . Also prove the converse.
\item
  An algebra of finite degree over a commutative field is Frobenius if and only if its regular representation is isomorphic to its coregular representation (use Exercise 3, c)).
\end{enumerate}

\begin{enumerate}
\def\labelenumi{\arabic{enumi})}
\setcounter{enumi}{4}
\tightlist
\item
  Let K be a commutative field and A be a K-algebra. We endow \(A^{*}\) with its canonical left A-module structure. For \(t \in A^{*}\) , we denote by \(\varphi_{t}\) the mapping \((x, y) \mapsto t(xy)\) from \(A \times A\) to K.
\end{enumerate}

\begin{enumerate}
\def\labelenumi{\alph{enumi})}
\item
  Prove that this defines a K-linear isomorphism from \(A^{*}\) to the set of K-bilinear mappings \(\varphi: A \times A \to K\) satisfying \(\varphi(xy, z) = \varphi(x, yz)\) for \(x, y, z \in A\) .
\item
  We say that the element \(t \in A^{*}\) is invertible if the homomorphism \(a \mapsto at\) from A to \(A^{*}\) is bijective. Let A be a K-algebra of finite degree. Prove that the following conditions are equivalent:
\end{enumerate}

\begin{enumerate}
\def\labelenumi{(\roman{enumi})}
\item
  A is Frobenius (Exercise 4).
\item
  A admits an invertible linear form.
\item
  A admits a linear form whose kernel does not contain any left ideals of \(\mathbf{A}\) .
\end{enumerate}

\begin{enumerate}
\def\labelenumi{\arabic{enumi})}
\setcounter{enumi}{5}
\tightlist
\item
  Let K be a commutative field. We say that a K-algebra A is symmetric if it admits a tracial invertible linear form (VIII, p.~115, Exercise 3).
\end{enumerate}

\begin{enumerate}
\def\labelenumi{\alph{enumi})}
\item
  Prove that a symmetric algebra has finite degree and is Frobenius.
\item
  Prove that the algebra of a finite group over \(\mathrm{K}\) is symmetric.
\item
  A semisimple algebra of finite degree over \(\mathrm{K}\) is symmetric.
\end{enumerate}

